\documentclass{beamer}
\usetheme{Warsaw}
\usepackage{caption}
\usepackage{tikz}
\usepackage{adjustbox}
\usetikzlibrary{arrows.meta, calc, quotes, tikzmark}
\usetikzlibrary{shapes.geometric, arrows, chains, positioning}

\setbeamerfont{caption}{series=\normalfont,size=\fontsize{7}{9}} 



\title[Springboard Capstone 2]{The Hunt for Exoplanets} 
\subtitle{Vetting Kepler Light Curves}

\author{Praveen Gowtham} % Your name
\institute[UCLA] % Your institution as it will appear on the bottom of every slide, may be shorthand to save space

\date{} % Date, can be changed to a custom date


\begin{document}

% Define block styles
\tikzstyle{decision} = [diamond, draw, fill=blue!20, 
text width=4.5em, text badly centered, node distance=3cm, inner sep=0pt]
\tikzstyle{block} = [rectangle, draw, fill=blue!20, 
text width=5em, text centered, rounded corners, minimum height=4em]
\tikzstyle{line} = [draw, -latex']
\tikzstyle{cloud} = [draw, ellipse,fill=red!20, node distance=3cm,
minimum height=2em]
	
	\begin{frame}
		\titlepage % Print the title page as the first slide
	\end{frame}

\begin{frame}{Motivations}
	\begin{figure}
		\includegraphics[scale=0.35]{pia17999_in_pj}
		\caption{An artist’s concept of Kepler-186f, the first validated Earth-size planet to orbit a distant star in the “habitable zone.”}
	\end{figure}


\end{frame}

\begin{frame}{Motivations}
	\begin{columns}
		\tikzset{box/.style={inner xsep=0pt}}
		\begin{column}{0.4\textwidth}
		
		  \begin{itemize}
			\item Search for habitable planets.
			\item The types and distribution of planets and planetary systems in Milky Way.
			\item Where does our solar system fit in?
		\end{itemize}	
		\end{column}
	\begin{column}{0.6\textwidth}
		\begin{figure}
			\includegraphics[scale=0.053]{planetseverywhere}
					\caption{Artist depiction of planetary system distribution in sector of Milky Way galaxy based off of exoplanetary data.}
		\end{figure}

	\end{column}
\end{columns}
\end{frame}

\begin{frame}{Detection technique: Light curve transit crossing}
	Light flux time series from star.  Periodic dips in light flux could be transiting exoplanet.
	
	\begin{columns}
		\begin{column}{0.5\textwidth}
			\begin{figure}
				\includegraphics[scale=0.11]{transit_illustration}
				\caption{Schematic of transit detection via space telescope photometry.}
			\end{figure}
		\end{column}
		\begin{column}{0.5\textwidth}
			\begin{figure}
				\includegraphics[scale=0.1]{exo_multiple_transit}
				\caption{Light curve with periodic transits from exoplanet HAT-P-7b.}
			\end{figure}
	\end{column}

		
	\end{columns}	
\end{frame}

\begin{frame}{Kepler mission (2009-2018)}
	\begin{columns}
		\begin{column}{0.4\textwidth}
			\begin{itemize}
				\item Kepler space telescope: dedicated photometry scanning a section of Milky Way.
				\item Photometry on $>$ 500,000 stars.
				\item $\approx$ 2,500 confirmed exoplanets by Kepler via light curve transit method.
			\end{itemize}	
		\end{column}	
		\begin{column}{0.6\textwidth}
			\begin{figure}
				\includegraphics[scale=0.09]{kepler_spacetelescope}
				\caption{Kepler Space Telescope.}
			\end{figure}
		\end{column}
			
		\end{columns}	
\end{frame}

\begin{frame}{The problem}
	\begin{columns}
		\begin{column}{0.5\textwidth}
			\begin{itemize}[<+-| alert@+>]
				\item Transit or noise? Some processing/statistics required.
				\item False positives (FPs).
				\begin{itemize}
					\item Secondary eclipse FPs
					\item Non-transiting Phenomena (NTP FPs)
				\end{itemize}
				\item Goal: Real exoplanets vs. different types of FPs
			\end{itemize}	
		\end{column}	
		\begin{column}{0.5\textwidth}
			\begin{figure}
						\includegraphics<1>[width=5.75cm]{exo_multiple_transit}
						\includegraphics<3>[width=5.75cm]{algol-curve}
						\visible<3>{\caption{Light curve from eclipsing binary star system (Algol A/B System)}}
						\includegraphics<4>[width=5.75cm]{phasefolded_ntp_ex2}	
						\visible<4>{\caption{Single oscillation cycle from pulsating star.}}		

			\end{figure}
		\end{column}
		
	\end{columns}	
\end{frame}

\begin{frame}{Kepler Data Validation Stream}
	Kepler mission uses an automated preprocessing and vetting pipeline.
\begin{itemize}
	\item Detect potentials: periodic statistically significant dip events. These are \textbf{Kepler Objects of Interest (KOIs)}.
	\item Basic preprocessing of KOI light curves.
\end{itemize} 
\begin{block}{Kepler pipeline}
\begin{adjustbox}{max totalsize={.9\textwidth}{.7\textheight},center}
	
    \begin{tikzpicture}[>=latex']
	
	\node [block] (start) {Kepler light curve \\ from pixel image};
	\node [block, right =2cm of start] (preprocess) {Whitening \\ + \\ detrend };
	\node [block, right =2cm of preprocess] (statdetector) {Periodic \\ event \\ detection};
	\node [block, right =2cm of statdetector] (otsu) {KOI catalog};
	\
	
	
	%% paths
	\path[draw,->]
	(start) edge (preprocess)
	(preprocess) edge (statdetector)
	(statdetector) edge (otsu)
	;
\end{tikzpicture}
		
\end{adjustbox}
\end{block}
\textbf{Focus on classification of KOIs}
\end{frame}	
	

\begin{frame}{Problem specification}
	\begin{block}{Label encoding + data breakdown}
		\begin{adjustbox}{max totalsize={.9\textwidth}{.7\textheight},center}
			
			\begin{tabular}{|c|c|c|c|}
			\hline
			Class name & Confirmed Planet  & Secondary Eclipse FP   &  NTP FP  \\
			\hline
			Target Label & 1  & 2 & 3    \\
			\hline
			Count & 2333 & 2140 & 880 \\
			\hline
			
		\end{tabular}
	
	\end{adjustbox}
	\end{block}
	Multi-label classification problem from light curves of KOIs.

\begin{itemize}
\item \textbf{Key challenge}: Light curve $\rightarrow$ processed light curve $\rightarrow$ relevant features
\item Train classifier on constructed features/label mapping
\end{itemize}

	
\end{frame}

\begin{frame}{Light curve processing: phase-folding + averaging}
\begin{columns}
	\begin{column}{0.5\textwidth}
	\begin{figure}
		\includegraphics[scale=0.35]{lc_detrend_secondaryFP_ex1}
		\caption{Whitened, detrended light curve.}
	\end{figure}
	\end{column}
	\begin{column}{0.5\textwidth}
	\begin{figure}
	\includegraphics[scale=0.35]{phasefolded_secondaryFP_ex1}
	\caption{Phase-folded and bin averaged (secondary eclipse FP).}
	\end{figure}
	\end{column}

\end{columns}	
\end{frame}

\begin{frame}{Constructed Features}
All features (except period) constructed from phase-folded, bin-averaged light curves. 
\vfill
\begin{adjustbox}{max totalsize={.9\textwidth}{.7\textheight},center}
	
		\begin{tabular}{|c|c|}
			\hline
			Period &  TCE period (days). \\
			\hline
			Duration & TCE duration (hours). \\
			\hline
			EOS & Even-odd statistic for secondary eclipse detection. \\
			\hline
			WSS & Weak secondary statistic  \\
			\hline
			min & Almost always corresponds to the depth of the primary transit.  \\
			\hline
			max & Maximum value in phase-folded/bin-averaged light curve. \\
			\hline
			LCBIN\_0 - LCBIN\_140 & 141 points of the xy-normalized primary transit close-ups  \\
			\hline
		\end{tabular}



\end{adjustbox}


\end{frame}

\begin{frame}{Weak Secondary Statistic (WSS)}
	\begin{columns}
		\begin{column}{0.5\textwidth}
		\begin{figure}
		\includegraphics[scale=0.35]{phasefolded_secondaryFP_ex1}
		\caption{Whitened, detrended light curve.}
	\end{figure}
	\end{column}
	\begin{column}{0.5\textwidth}
		\begin{figure}
	\includegraphics[scale=0.35]{secondarypeakonly_secondaryFP_ex1}
	\caption{Primary transit subtracted.}
\end{figure}
\end{column}
\end{columns}
\begin{block}{WSS Test}
	\begin{itemize}
		\item Subtract primary transit.
		\item Peak find and extract max amplitude peak (dip) + supports.
		\item Probability that amplitude (or more extreme) generated by Gaussian noise floor. 
	\end{itemize}
\end{block}
\end{frame}

\begin{frame}{Even Odd Statistic (EOS)}
Sometimes primary + secondary: similar amplitudes. Secondary at exactly half period. \\
\textbf{Problem: primary and secondary registered as same event by Kepler data validation.}
	\begin{columns}
	\begin{column}{0.5\textwidth}
		\begin{figure}
			\includegraphics[scale=0.35]{lc_detrend_secondaryFP_ex2}
			\caption{Primary and secondary transit amplitudes similar.}
		\end{figure}
	\end{column}
	\begin{column}{0.5\textwidth}
		\begin{figure}
			\includegraphics[scale=0.35]{evenoddstagger_secondaryFP_ex2}
			\caption{Alternate cycles binned into separate groups (even/odd). Each phase-folded, averaged.}
		\end{figure}
	\end{column}
\end{columns}
\end{frame}

\begin{frame}{Even Odd Statistic (EOS)}

	\begin{columns}
		\begin{column}{0.5\textwidth}
			\begin{figure}
				\includegraphics[scale=0.35]{lc_detrend_secondaryFP_ex2}
				\caption{Primary and secondary transit amplitudes similar.}
			\end{figure}
		\end{column}
		\begin{column}{0.5\textwidth}
			\begin{figure}
				\includegraphics[scale=0.35]{evenoddstagger_secondaryFP_ex2}
				\caption{Alternate cycles binned into separate groups (even/odd). Each phase-folded, averaged.}
			\end{figure}
		\end{column}
	\end{columns}
\begin{block}{EOS Test}
	\begin{itemize}
	\item Two sample t-test on even vs. odd phase transit depths.
	\item Extract p-value for significance of mean difference.
	\end{itemize}
	
\end{block}
\end{frame}

\begin{frame}{LCBIN Features}
	\begin{columns}
		\begin{column}{0.5\textwidth}
			\begin{itemize}
				\item Primary transit close-up: $\pm 2$ transit durations.
				\item XY-rescaled, resampled / binned to 141 points fixed length.
				\item \textbf{Shape difference:} U-shape vs. V-shape 
			\end{itemize}
		\end{column}
		\begin{column}{0.5\textwidth}
			\includegraphics[scale=0.375]{rescalexy_CP}
			\includegraphics[scale=0.375]{rescalexy_secondaryFP}
		\end{column}
	\end{columns}
	
\end{frame}

\begin{frame}{Feature Preprocessing Pipeline }
	
\begin{block}{Scikit-learn custom transformers + pipeline integration}

\begin{adjustbox}{max totalsize={.9\textwidth}{.7\textheight},center}
	
	\begin{tikzpicture}[auto]
		
		\node [block] (start) {WSS and EOS};
		\node [block, below = 1 cm of start] (start2) {LCBIN 0 to LCBIN 140};
		\node [block, below = 1 cm of start2] (start3) {Other features};
		\node [block, right =2cm of start] (preprocess1) {Log10 + Abs + Sqrt Transform};
				\node [block, right =2cm of start2] (preprocess2) {Dimensional reduction };
		\node [block, right =2cm of preprocess2] (normalizer) {Min/Max Normalizer};
		\node [block, right =2cm of normalizer] (output){Final Feature Matrix};
		
		
		
		%% paths
		\path[draw,->]
		(start) edge (preprocess1)
		(start2) edge (preprocess2)
		(preprocess2) edge (normalizer)
		(normalizer) edge (output)
		;
		
		\path [line] (preprocess1) -| (normalizer);
		\path [line] (start3) -| (normalizer);
	\end{tikzpicture}
	
\end{adjustbox}

\end{block}
\end{frame}


\begin{frame}{Statistical EDA: Secondary Eclipse Test Features}
\begin{columns}
\begin{column}{0.5\textwidth}		
	\begin{figure}
	\includegraphics[scale = 0.3]{psec_hist}	
	\end{figure}

	\begin{figure}
	\includegraphics[scale = 0.21]{psec_class_diff}			
	\end{figure}
\end{column}
\hspace{-20pt}
\vrule{}
\begin{column}{0.5\textwidth}
		\begin{figure}
		\includegraphics[scale = 0.3]{evenodd_hist}			
	\end{figure}
	\begin{figure}
	\includegraphics[scale = 0.21]{evenodd_class_diff.png}			
\end{figure}
\end{column}

\end{columns}

\end{frame}

\begin{frame}{Some other non-LCBIN features}
	\begin{columns}
		\begin{column}{0.5\textwidth}
			\begin{itemize}
				\item 'Period' good for class 3 (NTPs) selection.
				\item 'Min' feature / transit depth: exoplanets vs class 2 (secondary FPs)
				\item 'Duration', 'max': helps exoplanet vs. FP separation
			\end{itemize}
		
		\end{column}
		\begin{column}{0.5\textwidth}
	\begin{figure}
		\includegraphics[scale = 0.3]{ecdf_Period}			
	\end{figure}
	\begin{figure}
	\includegraphics[scale = 0.3]{ecdf_min}			
\end{figure}
\end{column}
	\end{columns}
	
\end{frame}


\begin{frame}{Locality Preserving Projections (LPP)}
	\begin{block}{LPP: Dimensional reduction technique}
		Map to low dimension maintains closeness of objects that are close in high dimensions. Reduce 141-D LCBIN to 2D space.
	\end{block}
\begin{columns}
	\begin{column}{0.5\textwidth}
		\includegraphics[scale=0.25]{rescalexy_CP} \tikzmark{b}
		\includegraphics[scale=0.25]{rescalexy_secondaryFP} 
	\end{column} \tikzmark{f}
	\begin{column}{0.5\textwidth}  
		\includegraphics[scale=0.3]{lppjointplot}

	\end{column}  
\end{columns}

    \begin{tikzpicture}[overlay,remember picture]
	\draw[very thick, -Stealth]         ($({pic cs:b})+(1ex,1ex)$)
	to [bend left, "LPP"]  ($({pic cs:f})+(1ex,+1em)$);
\end{tikzpicture}
\end{frame}

\begin{frame}{Class separation in LPP space}
	\begin{figure}
	\includegraphics[scale=0.3]{lpp_hexbinplot}
	\caption{Hexbin count distribution for each class in 2D LPP space.}
	
	\end{figure}
Distinction between class 1 (exoplanets) and FPs. Outliers dominated by NTP FPs.
\end{frame}


\begin{frame}{Modeling Pipeline}
\begin{block}{Hyperparameter tuning: CV grid search}	
	\begin{itemize}
		\item 75-15 train/hold-out set. 
		\item Scikit-learn full pipeline:
		
			\begin{adjustbox}{max totalsize={.9\textwidth}{.7\textheight},{margin= 3mm}, center}
			
			\begin{tikzpicture}[auto]
				
				
				\node [block] (start) {Pre-process pipeline};
				\node [block, right =2cm of start] (model) {Classifier model};
				\
				
				
				%% paths
				\path[draw,->]
				(start) edge (model)
				
				;
			\end{tikzpicture}
			
		\end{adjustbox}
	\item 5-fold cross validation: 75-15 splits on train set. 
	\item Eval. metric: f1-score micro-averaged across classes.
	\item Pipeline + CV fit: avoids data leakage.
		
	\end{itemize}
	
\end{block}

\begin{block}{Classifier Models + Hyperparameters}
\textbf{Random Forest}: max\_depth, max\_features, n\_estimators \\ 
\textbf{RBF-Kernelized Soft-Margin SVM}: C, gamma \\ 
\textbf{XGBoost (stumps)}:  learning\_rate, n\_estimators
\end{block}

\end{frame}

\begin{frame}{Best Performing Model}
		\begin{itemize}
			\item Final evaluation of tuned models on hold-out set.
			\item XGBoost using decision stumps performs the best.
			\item Optimal hyperparameters:
			\begin{itemize}
				\item Learning rate: 0.1
				\item Number of estimators: 1000
			\end{itemize}
		\end{itemize}
	
	\begin{block}{XGBoost performance on hold-out}
		\centering
		\begin{tabular}{|l|r|r|r|r|}
			\hline 
			class &  precision &    recall &  f1-score &    support \\ \hline
			
			1            &   0.943820 &  0.957265 &  0.950495 &  351.00000 \\ \hline
			2            &   0.867868 &  0.900312 &  0.883792 &  321.00000 \\ \hline
			3            &   0.773913 &  0.674242 &  0.720648 &  132.00000 \\ \hline
 
			
		\end{tabular}
	\end{block}

\end{frame}
\begin{frame}{Performance Summary + Feature Evaluation}
	\begin{columns}
		\begin{column}{0.5\textwidth}
			\begin{itemize}
				\item Exoplanet vs. FP: excellent.
				\item Secondary eclipse FP: pretty good.
				\item NTP FPs: some confusion w/ secondary FPs.
			\end{itemize}
		\end{column}

		\begin{column}{0.5\textwidth}
			\begin{figure}
					\includegraphics[scale = 0.3]{exo_xgb_cfmat}
			\end{figure}

		\end{column}
	
	\end{columns}
\begin{columns}
	\begin{column}{0.5\textwidth}
		\begin{itemize}
			\item Feature importances on similar scale.
			\item Different tree ensemble methods: ordering robust.
		\end{itemize}
	\end{column}
	\begin{column}{0.5\textwidth}
	\begin{figure}
		\includegraphics<1>[scale = 0.25]{xgb_feature_imp}
		\includegraphics<2>[scale = 0.25]{randforest_feature_imp}
	\end{figure}
	\end{column}
\end{columns}

\end{frame}

\begin{frame}{Future Directions}
	\begin{columns}
	\begin{column}{0.5\textwidth}
	\begin{itemize}
		\item Start vetting from raw data. End-to-end pipeline.
		\item New missions, similar mission data validation streams.
		\item Model adaptation to new missions.

	\end{itemize}	
	\end{column}
	\begin{column}{0.5\textwidth}
		\begin{figure}
		\includegraphics[scale = 0.5]{TESS}
		\caption{Transiting Exoplanet Survey Satellite (TESS)}
	\end{figure}
	\end{column}
\end{columns}
\end{frame}
\end{document}